\chapter{Criteria and thresholds}
\label{ch:criteria}

Every number the program judges against, with its source class. \textsf{[Manual]}
= a program manual (section and URL are printed in the report);
\textsf{[Literature]} = a DOI-backed reference (the complete citation and
BibTeX entry are printed in the report); \textsf{[Measured]} = a record of this
project (a fixture or an internal record, cited by path). Thresholds marked
\emph{provisional} are this project's working values, to be re-calibrated when
evidence accumulates.

\section{Occupations and active spaces}

\begin{center}
\begin{tabularx}{\textwidth}{@{}>{\raggedright\arraybackslash}p{3.2cm} >{\raggedright\arraybackslash}p{2.0cm} X >{\raggedright\arraybackslash}p{3.2cm}@{}}
\toprule
Quantity & Value & Meaning / limits & Source \\
\midrule
Active-orbital window & 0.02--1.98 & ORCA's auto-ICE convention, borrowed here
  as the definition of ``fractionally occupied''; an active orbital outside it
  is reported as nearly empty / nearly full (a ``please confirm'' warning, not
  an error -- a full f shell in the window is deliberate in magnetic work) &
  \textsf{[Manual]} ORCA 6.1 \S3.14 \\
f-electron count & filling rule & Standard filling for the f$^n$ configuration
  of the trivalent (or stated-valence) ion; drives the active-space suggestion
  only & \textsf{[Measured]} module tests \\
Cross-structure mapping & $\tau$ read from the curve (the source's examples:
  $\tau \le 0.5$) & Menu~\menu{17}: orbitals of different structures match when
  their kinetic energies and their shell-wise populations agree within one
  threshold $\tau$; the tau curve's plateaus are the feature, not one number.
  The populations here are shell-wise L\"owdin (the source's IAO choice is
  licensed by its own text as ``any orbital-wise population analysis''), so
  $\tau$ is data-calibrated here, never copied from the source's absolute
  values & \textsf{[Literature]} Bensberg \& Reiher 2023, Eqs.~(1)/(8) \\
Non-matchable block & one class together & Menu~\menu{17}: the orbitals that
  change along the path cannot be matched in every structure; they join the
  active space as one block whenever one of them is selected anywhere (the
  source's union rule).  Measured on the N$_2$ scan: cores map, the
  one-centre hybrids are the block & \textsf{[Literature]} same;
  \textsf{[Measured]} the frozen scan fixtures \\
Active-space overlap & $|\det S_{\mathrm{act}}|$ near 1 preserved, near 0 an
  exchange & Menu~\menu{17}'s optional \code{active} block (the second source's
  Supporting Information): the active block of two \emph{adjacent} structures
  matched by overlap, the current geometry's $S$ used for the whole matrix (its
  small-step approximation).  The printed bands are this project's reading of
  the source's qualitative scale -- measured: cores 1.0000 for every pair, the
  bond triad 0.9955 across 0.01~\AA{} and 0.72 across 0.5~\AA{}; a synthetic
  active$\leftrightarrow$inactive exchange lands at $6\times10^{-5}$ &
  \textsf{[Literature]} Paz et al.\ 2021 SI Eq.~(2), main text pending;
  \textsf{[Measured]} the frozen scan \\
Guess transfer (WASP) & no threshold; $1/d$ weights & Menu~\menu{18}: the
  target's guess as the 1/d-weighted interpolation of the neighbours over
  $\delta$, orthonormalised in the target's metric; the written mkl converts
  to a \code{.gbw} ORCA reads as \code{INITIAL GUESS: MOREAD} (round trip
  measured at $5\times10^{-8}$).  \textbf{A guess steers the basin}: on the
  frozen example it converged a \emph{second} RHF solution 19~mEh below the
  scan's branch -- branch continuity must be checked after any transfer &
  \textsf{[Literature]} the review's WASP (Eqs.~(11)--(13));
  \textsf{[Measured]} the frozen chain \\
Dipole-moment selection (DM-AS) & no threshold; smallest deviation wins
  (ties: smaller space) & Menus~\menu{19}/\menu{20}: candidates (the source's
  PASS family, 6--14 electrons with $n_\mathrm{e}/2+3 \le n_\mathrm{o} \le
  14$) ranked by $|\mu - \mu_\mathrm{ref}|$ (directional: projected vector
  difference); the reference's origin (experiment / functional) must be
  stated.  Measured on the H$_2$O fixture chain: PBE0 reference 2.0801~D and
  the winner $(6e, 8o)$ at 0.0081~D.  The per-state variants (EDM/D$^2$DM)
  are not implemented: ORCA prints only the state-average dipole for
  state-averaged runs and no TDDFT state dipoles (measured) &
  \textsf{[Literature]} Kaufold et al.\ 2023, 2026; \textsf{[Measured]} the
  H$_2$O chain \\
APC ranked-orbital selection & no threshold; CSF cap & Menu~\menu{21}:
  candidate orbitals (all doubly occupied plus the lowest virtuals, default
  23, the source's general scheme) ranked by the approximate pair coefficient
  $c_{ia} = -x_a/(D_{ia}+\sqrt{x_a^2+D_{ia}^2})$, $x_a = \tfrac12 K_{aa}$,
  $D_{ia} = \varepsilon_a-\varepsilon_i$; the selection drops the lowest until
  the CSF count $\binom{L}{\alpha}\binom{L}{\beta} -
  \binom{L}{\alpha+1}\binom{L}{\beta-1}$ fits the cap, keeping at least one
  occupied and two unoccupied orbitals (the source's stability floor).
  The source's error record travels with every report (doubly-occupied
  entropy $R^2=0.64$, MAE $0.0240$ against DMRG; precision about $88\%$);
  entropies agreeing to $10^{-12}$ rank by orbital index (BLAS determinism).
  Measured anchor: the N$_2$ max(10,10) selection lands exactly on the
  source's $(10,10)=19404$ entry &
  \textsf{[Literature]} King \& Gagliardi 2021 \S3.2, 3.12;
  \textsf{[Measured]} \code{tests/test\_apc.py} \\
ASS1ST quasi-NOON band & $[T, 2-T]$; $T=0.05$ conservative, $0.03$ relaxed
  (two independent lines allowed, e.g. $0.03/1.96$) & Menus~\menu{22}/\menu{23}:
  block eigenvalues of the perturbation-theory density are taken as strongly
  correlated inside the band; bookkeeping adds $+2$e$+1$o (internal) or
  $+0$e$+1$o (external) and reassigns drifting active orbitals out
  ($-2$e$-1$o / $-0$e$-1$o).  Substitution carried: ORCA answers the
  unrelaxed density only for FIC-NEVPT2 (measured), the source uses
  SC-NEVPT2 --- same object family, different contraction, counts near the
  thresholds may shift.  Measured anchors: N2 CAS(6,6) internal block
  2.00000, 2.00000, 1.98196, 1.97680 and the $(6,6)\to(4,4)$ suggestion,
  from both the single-root and the three-singlet round; the CASSCF
  reference density pins the sidecar convention $D_{AO}=S D_{\mathrm{json}} S$
  (exactly diagonal in its natural-orbital basis, couplings $10^{-15}$) &
  \textsf{[Literature]} Khedkar \& Roemelt 2019 \S2.4, 2020 \S2.4;
  \textsf{[Measured]} \code{tests/test\_ass1st.py} \\
\bottomrule
\end{tabularx}
\end{center}

\section{Entropy and multi-reference character}

\begin{center}
\begin{tabularx}{\textwidth}{@{}>{\raggedright\arraybackslash}p{3.2cm} >{\raggedright\arraybackslash}p{2.0cm} X >{\raggedright\arraybackslash}p{3.2cm}@{}}
\toprule
Quantity & Value & Meaning / limits & Source \\
\midrule
Single-orbital entropy line & max $s_i > 0.14$ & For the \emph{true four-state}
  entropy $s_i = -\sum_j \rho_{jj} \ln \rho_{jj}$ over the four occupation
  states of an orbital (ceiling $\ln 4$). An ORCA output gives only
  spin-summed occupations, which do not determine that entropy; menu~\menu{1}
  computes a provable upper bound instead and applies 0.14 one-sidedly --
  bound $\le 0.14$ excludes, bound $> 0.14$ is indecisive &
  \textsf{[Literature]} Stein \& Reiher 2016; Wardzala et al.\ 2026 Eq.~(6) \\
Exact four-state entropy & $s_i$ computed directly & Menu~\menu{12}
  reconstructs $s_i$ from the user's FCIDUMP (converged CASSCF plus the
  \code{!FCIDUMP} dump run) and reads 0.14 against the \emph{localized}-basis
  spectrum; the reconstruction is gated by engine cross-checks (the CI energy
  must reproduce the printed CASSCF energy, the natural occupations the
  \code{N(occ)} line) &
  \textsf{[Measured]} route validation 2026-09-26 (energy agreement $4\times
  10^{-13}$ on the N2 fixture; \code{tests/test\_entropy\_rdm.py}) \\
Relativistic tier & tier table (menu~\menu{3}) & The relativistic
  Hamiltonian the need implies: one-electron variants for scalar work;
  X2Ccorr for light main-group spin-orbit (about a quarter of the O$_2$
  splitting); X2CAMF for the f block, where X2Ccorr moves Nd\textsuperscript{3+}
  splittings by only $1.6$--$4.9$~cm$^{-1}$ (over-configuration); X2CMP + QED
  for the highest accuracy and for core spectroscopy; amfX2C ruled out
  (numerically unstable scalar part) &
  \textsf{[Literature]} Wang et al., arXiv:2602.24236v1; \textsf{[Measured]}
  the Nd\textsuperscript{3+} template's 52-state window computed from the Hund
  term \\
Orbital bonding label & pair cutoff $6$~\AA; single-atom share $95\%$;
  cross band $0.01$; cancellation limit $5$ & Menu~\menu{15}'s per-orbital
  bonding label: the atom-pair cross population decides bonding from
  antibonding unless one atom carries more than $95\%$ of the L\"owdin
  population.  The band and the limit are this project's additions, sized on
  the N2 fixture (a symmetric core measures $3\times10^{-4}$; the orbitals run
  up to an absolute block sum of $2.2$ while one node-heavy virtual measures
  $158$) &
  \textsf{[Literature]} Osaro et al., arXiv:2606.07879v1 \S3.1;
  \textsf{[Measured]} the N2 export \\
Activity size vs error & not monotone & The descriptor-panel source's own
  benchmark: more orbitals is not automatically better (CH$_4$: 17 orbitals at
  twice the error of 8; NH$_3$: 13 at fourteen times the error of 6), so a
  diagnosis must not treat ``smaller than expected'' as a defect by itself &
  \textsf{[Literature]} Osaro et al., arXiv:2606.07879v1 \S5 \\
Diffuse candidate (rule G) & no threshold (ranking only) & The diffusest
  Gaussian exponent of an active orbital's dominant (element, angular momentum)
  channel, ranked across the active orbitals (menu~\menu{1}, section A8). The
  source asks for identification and sorting and gives no line; the exponent's
  absolute value does not transfer between elements (measured: $0.19$ for N in
  def2-SVP against $0.022$ for a Eu valence s in SARC2) &
  \textsf{[Literature]} Lee \& Rondinelli arXiv:2609.13357 (rule G);
  \textsf{[Measured]} the two basis fixtures \\
AVAS truncation & $0.1$ (default; the source's range $0.05$--$0.1$) &
  Overlap eigenvalues of the target-shell projection (menu~\menu{14}); orbitals
  below the threshold leave the active space. The source notes the threshold is
  the only user-chosen number besides the target and the SCF state, and that a
  cut of low-overlap orbitals is what keeps CASCI below the variational HF
  energy &
  \textsf{[Literature]} Sayfutyarova et al.\ 2017 \S2.4; range $0.05$--$0.1$ \\
AVAS virtual side & no published line & A virtual-side spectrum with no
  eigenvalue above the threshold means the antibonding target character is
  ligand-centred (the source's $[\mathrm{CuCl}_4]^{2-}$ case): add ligand AOs
  to the target or record the target as unsuited &
  \textsf{[Literature]} Sayfutyarova et al.\ 2017 \S4.5 (worked examples) \\
Environment spin entropy & $0.05$ (provisional);
  source anchors $0.007$ / $2.766$ & $\Delta S_E$ of Eq.~\eqref{eq:delta-se}
  (menu~\menu{12}): zero exactly when the environment's two spin densities are
  equal. The provisional line sits an order of magnitude above the source's
  correct-solution value and eight orders above the exact route's noise; the
  zero is exact and a closed-shell state gives it in every basis &
  \textsf{[Literature + provisional]} Ai et al.\ 2025 Eq.~(9); anchor values
  measured there \\
Subspace fraction & $1 - \sigma_F$ below $10^{-4}$ contained;
  $10^{-4}$--$10^{-2}$ approximate & $\sigma_F$ (menu~\menu{13}) is a ranking
  quantity: the source reports that the ordering of $1 - \sigma_F$ tracks the
  ordering of the energy error ($2.9\times10^{-3}$ loose start,
  $6.8\times10^{-5}$ after a localisation). Not an absolute pass mark, and
  blind to a space being too large &
  \textsf{[Literature + provisional]} Guan \& Jiang Eq.~(10);
  bands provisional \\
Space-change singular value & $0.9$ / $0.5$ (provisional) &
  Smallest singular value of $C_{\mathrm{final}}^{\mathsf{T}} S\,
  C_{\mathrm{initial}}$ (menu~\menu{13}) for equal-size sets: near $1$ the
  space did not move, near $0$ an initial active orbital was replaced. The
  source's own runs measure $0.65$--$0.99$ &
  \textsf{[Literature + provisional]} Sayfutyarova et al.\ 2017 Eq.~(16) \\
Atomic-term $f$ character & $0.9$ (provisional) & Mean projected $l$-shell
  character of the active orbitals (menu~\menu{12}): below the line the Hund
  verdict is withheld, because the quantum numbers then describe only the
  projected part. The $f$-block wrong-solution mode lands here with a
  character near $0$ &
  \textsf{[Measured + provisional]} Eu$^{3+}$ d-solution fixture measured at
  $0.000$; clean-shell cases at $1.000$ \\
Effective $J$ & $|J_{\mathrm{eff}} - (L + S)|$ and
  $|J_{\mathrm{eff}} - |L-S||$ within $0.1$; $L$, $S$ within $0.1$ of Hund &
  Menu~\menu{12}. Only a mismatch in $L$ or $S$ is a defect; a non-stretched
  $M_J$ component has no definite $J$ in a scalar calculation and is reported
  as such. The source's discriminating case is the $L$ drift ($5.1 \to 4.6$)
  &
  \textsf{[Literature]} Peng et al.\ 2025 Table 1 (Dy\textsuperscript{3+}:
  $J = 7.5$ only for the HF density) \\
Weak-correlation cut & 1--2\% of max & Relative cut used with the ranking
  proxy when no plateau is found & \textsf{[Literature]} Stein \& Reiher 2016 \\
$Z_{s(1)}$ line & 0.2 & A \emph{different} diagnostic; listed here only to say
  it is never mixed with the 0.14 line & \textsf{[Literature]} Stein \& Reiher
  2017 \\
T1 screening & 0.02 & Closed-shell screening convention for single-reference
  adequacy (CCSD singles norm / $\sqrt{N_{\text{corr}}}$); a screening result,
  not a law. ORCA prints T1, not D1 & \textsf{[Literature]} Lee \& Taylor 1989 \\
Fractional-orbital count & no published line & Reported as a count; an
  occupation within $10^{-3}$ of an integer counts as integer-valued (single
  configuration) & \textsf{[Measured]} module convention \\
\bottomrule
\end{tabularx}
\end{center}

\section{CASSCF, NEVPT2 and CASPT2}

\begin{center}
\begin{tabularx}{\textwidth}{@{}>{\raggedright\arraybackslash}p{3.2cm} >{\raggedright\arraybackslash}p{2.0cm} X >{\raggedright\arraybackslash}p{3.2cm}@{}}
\toprule
Quantity & Value & Meaning / limits & Source \\
\midrule
CASSCF convergence & energy \emph{or} gradient marker & ORCA prints two
  different convergence markers; energy convergence does not imply orbital
  convergence, and the perturbation layer depends on the orbitals (order of
  kcal/mol) & \textsf{[Measured]} group records + fixtures \\
CASPT2 reference weight & recommend $> 0.9$ & Below 0.9 the state is not
  dominated by the reference; ORCA does \emph{not} abort on intruders, so the
  check is mandatory & \textsf{[Manual]} ORCA 6.1 (CASPT2) \\
CASPT2 smallest denominator & alert below 0.01 & A small denominator is the
  intruder-state risk the manual describes; the program alerts, it does not
  decide & \textsf{[Manual]} same section \\
NEVPT2 hole/particle classes & alert when $> 0$ & Positive V1\_i (1-hole) or
  V m1\_a (1-particle) class contributions indicate a (false) intruder state &
  \textsf{[Manual]} ORCA 6.1 \S3.16.6 \\
Embedding error (DMET, 4f SIMs) & CASSCF-SO 0.6--7.8; NEVPT2 8.8--62.7~cm$^{-1}$
  & Against all-electron on three real Dy/Er systems. NEVPT2's embedding error
  grows systematically --- dynamic correlation is less local than the static
  part --- so the metal-only cluster is not enough: expanding the cluster to the
  metal plus its nearest coordinating atoms brings 3Dy back $62.7 \to
  13.2$~cm$^{-1}$ (4.75$\times$), the default whenever NEVPT2 is used. A wrong
  low-level SCF solution costs 42$\times$ on the fitted parameters (357.7
  against 8.6~cm$^{-1}$) & \textsf{[Literature]} Ai et al.\ 2025 (Tables 1--3) \\
\bottomrule
\end{tabularx}
\end{center}

\section{SCF}

\begin{center}
\begin{tabularx}{\textwidth}{@{}>{\raggedright\arraybackslash}p{3.2cm} >{\raggedright\arraybackslash}p{2.0cm} X >{\raggedright\arraybackslash}p{3.2cm}@{}}
\toprule
Quantity & Value & Meaning / limits & Source \\
\midrule
Suspected pseudo-convergence & converged within $\le 3$ cycles & Counter-check
  the cycle count; a real transition-metal case converged in 2 cycles &
  \textsf{[Measured]} SCINE capability report \\
Rescue paths & see Chapter~\ref{ch:menu9} & Better starting orbitals first;
  \code{SlowConv} with its caution; \code{MaxIter} alone is refused (the
  manual: ``will not help in many cases'') & \textsf{[Manual]} ORCA 6.1
  \S2.6.x; \textsf{[Measured]} fixtures \\
\bottomrule
\end{tabularx}
\end{center}

\section{Frequencies and transition states}

\begin{center}
\begin{tabularx}{\textwidth}{@{}>{\raggedright\arraybackslash}p{3.2cm} >{\raggedright\arraybackslash}p{2.0cm} X >{\raggedright\arraybackslash}p{3.2cm}@{}}
\toprule
Quantity & Value & Meaning / limits & Source \\
\midrule
Imaginary-mode reading & 0 / 1 / $\ge 2$ & 0 = minimum; exactly 1 = first-order
  saddle (a TS candidate); $\ge 2$ = higher-order saddle. The count must come
  from the last frequency block (an OptTS job prints the initial Hessian block
  first) & \textsf{[Manual]} ORCA 6.1 TS tutorial \\
Small-imaginary warning & $|\nu| < 50$~cm$^{-1}$ & \emph{Provisional}: in this
  regime a numerical Hessian can report imaginary modes whose energy curvature
  is positive (measured: spurious modes between about $-27$ and
  $-278$~cm$^{-1}$ on a loose two-fragment system, settled by an
  energy-curvature scan). The rule says ``verify'', never ``wrong'' &
  \textsf{[Measured]} group audit report \\
TS verification & frequency after the optimisation & A converged TS optimisation
  is not evidence by itself; the recommended route uses an exact starting
  Hessian, and adding \code{Freq} verifies exactly one imaginary mode &
  \textsf{[Manual]} ORCA 6.1 TS tutorial \\
Barrier reference point & always state it & The same barrier data read
  $+5.90$ or $-22.14$~kcal/mol depending on whether the reference is a
  pre-complex or the separated reactants & \textsf{[Measured]} SCINE
  capability report \\
\bottomrule
\end{tabularx}
\end{center}

\section{Crystal fields}

\begin{center}
\begin{tabularx}{\textwidth}{@{}>{\raggedright\arraybackslash}p{3.2cm} >{\raggedright\arraybackslash}p{2.0cm} X >{\raggedright\arraybackslash}p{3.2cm}@{}}
\toprule
Quantity & Value & Meaning / limits & Source \\
\midrule
Allowed $B_k^q$ & $C_3$: 9; $O_h$: 4; none: 27 & Even $k$ only (time reversal),
  $k \le 6$ (4f) & \textsf{[Literature]} Peng et al.\ 2025; Chilton 2025 \\
Cross-scheme comparison & $(2,0)$ only, roughly & Same wavefunction, different
  projection: $B_2^2$ moved $+198 \to -43$~cm$^{-1}$; $B_2^0$ moved 3\%. The
  licence is exactly the leading axial parameter & \textsf{[Literature]}
  Chilton 2025 (Table S1) \\
Ill-conditioning limit & cond $> 10^8$ & Flagged, with the rank-deficiency note
  for eigenstate-only sampling & \textsf{[Measured]} module tests \\
Point-charge estimate & qualitative & Measured point-charge/ab initio
  deviations make the model unusable for quantitative parameters; $\langle
  r^k\rangle$ is user-supplied & \textsf{[Literature]} Ungur \& Chibotaru 2017;
  Scheie 2021 \\
$g_T$ criterion (doublet ladder) & line $g_T\theta_3 = 20$ &
  $g_T = (g_1 + g_2 + g_3 \sin\theta_3) / 3$ with $\theta_3$ in
  \emph{degrees}; below the line the doublet can still act as a step of the
  barrier, at or above it quantum tunnelling is opened. Empirical, not derived
  (the source states the $g_1,g_2$ plane-out correction is not available), and
  the line is not a universal constant (the internal dipolar field varies with
  packing); calibrated on 19 mononuclear Dy(III) systems, where the line sits in
  the measured gap $[15.31, 30.25]$ & \textsf{[Literature]} Chilton 2025
  (Table S2) \\
Doublet-table regression (menu 16) & 38/38 tabulated $g_T$ reproduced & The
  source's whole calibration table is the fixture: every tabulated $g_T$ is
  recomputed from the raw $g$/$ \theta_3$ columns (largest disagreement
  $8\times10^{-4}$), the gap is reproduced, and the one doublet on the wrong
  side of the line is the source's recorded outlier
  ([Dy(Cp\textsuperscript{ttt})\textsubscript{2}]\textsuperscript{+}, product
  2.52) & \textsf{[Measured]} regression check 2026-09-27 \\
\bottomrule
\end{tabularx}
\end{center}

\section{Cross-level consistency}

\begin{center}
\begin{tabularx}{\textwidth}{@{}>{\raggedright\arraybackslash}p{3.2cm} >{\raggedright\arraybackslash}p{2.0cm} X >{\raggedright\arraybackslash}p{3.2cm}@{}}
\toprule
Quantity & Value & Meaning / limits & Source \\
\midrule
Top-$N$ dropout & $N = 3$ & A solution in the cheap level's top 3 that leaves
  the expensive level's top 3 is reported; so is a disagreement about the best
  solution. In the source case the cheap-level best was not the expensive-level
  best, and the cheap-level second-best fell to position 7 &
  \textsf{[Literature]} Lu et al.\ 2025 (SI Table S18) \\
\bottomrule
\end{tabularx}
\end{center}

\section{Local spin analysis}

\begin{center}
\begin{tabularx}{\textwidth}{@{}>{\raggedright\arraybackslash}p{3.2cm} >{\raggedright\arraybackslash}p{2.0cm} X >{\raggedright\arraybackslash}p{3.2cm}@{}}
\toprule
Quantity & Value & Meaning / limits & Source \\
\midrule
Block selection & last block & SCF jobs print several blocks (guess, after SCF,
  final); only the last describes the converged wavefunction & \textsf{[Measured]} fixtures \\
Spin purity & $\langle S^2 \rangle$ vs $S(S+1)$ & Read from the printed numbers;
  a deviation is a property of the method (broken symmetry is deliberately not
  spin-pure), not a file defect & \textsf{[Measured]} fixture values \\
Metal spin share & own construction, provisional & the metal share
  $(\sum_B \langle S_M S_B \rangle) / \langle S^2 \rangle$; undefined (0/0) for
  a total singlet, where the section reports the $S_{\mathrm{eff}}$ magnitudes
  instead & this manual's program, declared in its output \\
$\Delta S_E$ boundary & not available from tables & Local spin analysis is not
  the environment spin-polarisation entropy; the real quantity needs the
  two-step \code{orca\_2json} + \code{orca\_loc} route; the measured scale it
  stands in for: $\Delta S_E$ 2.766 (wrong solution) vs 0.007 (correct), and a
  downstream MAE of 357.7 vs 8.6~cm$^{-1}$ & \textsf{[Literature]} Ai et al.\\
  2025; \textsf{[Manual]} ORCA 6.1 \S5.1.10 \\
\bottomrule
\end{tabularx}
\end{center}

\section{Composition and analysis conventions}

\begin{center}
\begin{tabularx}{\textwidth}{@{}>{\raggedright\arraybackslash}p{3.2cm} >{\raggedright\arraybackslash}p{2.0cm} X >{\raggedright\arraybackslash}p{3.2cm}@{}}
\toprule
Quantity & Value & Meaning / limits & Source \\
\midrule
Shell assignment & largest (atom, shell) weight & Rank within an occupation
  partition only -- ranking the whole space by weight picks up virtual orbitals
  & \textsf{[Measured]} group's EuF record \\
Composition table & print flag needed & \code{\%output Print[P\_ReducedOrbPopMO\_L] 1};
  without it the A1 section is skipped with an explanation & \textsf{[Manual]}
  ORCA 6.1 \S2.9 \\
\bottomrule
\end{tabularx}
\end{center}
